\section{Conclusion}\label{sec:conclusion}

When two models are compared on a benchmark whose labels are partly determined by a feature
both of them observe, the reported gain answers a question nobody asked. It sums a closer fit
to the label frequencies within groups of that feature and a closer match between predictions
and the individual cases they were made for --- different achievements with different
consequences, since the first is skill a frequency table also supplies and that a shift in
label frequencies removes, and the second is not.

This paper separates them exactly. Scoring each model's prediction for one test case against
the label of a different case in the same group leaves the group and its label frequencies
intact while breaking the correspondence between prediction and case, so the gain surviving
that relabelling is the transported component and the remainder is a within-group covariance
gain. The identity holds in the observed sample for any proper score, with no fitted
reference model, no held-out fold and no tuning parameter, at $O(NK)$ on cached predictions.
The remainder is an order-two U-statistic equal, for every Bregman score, to a within-group
covariance between the models' probability change and the label indicator, and it comes with
a one-pass influence function, cluster-robust intervals, a randomization test and a limit
theorem.

\paragraph{What it buys} Pointed at two released MOTIFS checkpoints with their authors' own
code, the decomposition shows a ten-point mean-recall improvement corresponding to a
proper-score difference that is overwhelmingly transported. On predictors whose information
content we control it returns exactly zero covariance gain when both models are functions of
the grouping variable alone, and shows a frozen-CLIP model that loses on every aggregate
measure discriminating individual cases significantly better than the model it appears to
lose to.

\paragraph{Weaknesses} Three are structural. The estimator is not invariant to
recalibration, so we compare only confidence-matched pairs and released checkpoints cannot be
audited on raw outputs. It credits any signal varying inside a cell, shortcuts included, and
our own sensitivity bound is not reassuring: an unrecorded covariate correlated at $0.061$
would account for the entire covariance gain we report on Visual Genome. And the estimand is
not determined by the granularity of $\phi$ --- two equally fine, fully identified partitions
of the same four cases can disagree completely --- so declaring $\phi$ is a modelling decision
that choosing the finest defensible partition does not settle.

\paragraph{Future work} A rank-based within-cell contrast would be invariant to the
recalibration this construction is sensitive to, at the cost of the exact additivity that is
the point here; reporting both is what we would recommend to a benchmark able to choose.
Approximate matching would extend the identity to continuous grouping variables, trading
exactness for modelling assumptions. And porting multicalibration-style guarantees to
pairwise attribution would address the choice of $\phi$ more convincingly than a declaration
rule can. The obstacle to all of it is the one that limits the method today: benchmarks
archive ranked predictions, not probabilities.
